% \ad{Review}
% We studied semantic invariance in language model representations from a local geometric perspective. We defined invariant latent features as representations that are stable under semantic-preserving transformations while remaining discriminative across meanings, and we identified them with a contrastive sensitivity decomposition based on a regularized generalized eigenvalue problem. Across nine instruction-tuned models, invariant structure was strongest in mid-to-late layers, semantic displacement lay largely outside the nuisance subspace, and interventions on the invariant component changed model outputs in a way consistent with a causal role. Invariant representations also supported zero-shot model attribution, including for fine-tuned and distilled variants. These findings describe a geometric characterization of semantic invariance. Future work includes multilingual, long-context, and naturally occurring perturbations; larger models; controlled comparisons within a model family and between base and instruction-tuned models; and a more complete disentanglement of the nuisance representation.

We presented a local geometric framework for identifying semantic-invariant features in language model representations through contrastive sensitivity decomposition. Experiments across nine language models demonstrate geometric separation between semantic and nuisance variation, causal sensitivity to semantic changes, and robust zero-shot model attribution. Together, these findings support a local geometric characterization of semantic invariance and its model-specific structure.
